% GENERATED FILE. Edit canonical lessons, then run npm run generate:books.
% canonical-source-hash: fnv1a64:0c25524db9096c04
% canonical-lessons: BN-W01-na, BN-W01-aa-matra, BN-W01-aa, BN-W01-ha, BN-W01-ma, BN-W01-sa, BN-W01-ka, BN-W01-hasanta, BN-W01-ra, BN-W01-nomoshkar-read

\chapter{Hanging From the Head-Line, One Piece at a Time}
\label{ch:bn-script-first}
\hlchaptermodality{16}

\begin{chapteropening}
\textbf{By the end of this chapter:} \emph{I can write nine Bengali pieces from memory — six consonants, a vowel sign, its independent twin and the hasanta — and read the words for no and the greeting by assembling them from those pieces, including the greeting's conjunct.}

The greeting assembled from seven pieces the reader can each write from memory, including the conjunct and the letter written s but said sh, and introducing nothing new.
\end{chapteropening}

\section[nɔ]{\bn{ন} — the consonant na}

\label{lesson:BN-W01-na}



\begin{quote}
Nothing is assumed but the words you can already say. This chapter turns them into words you can read.
\end{quote}

\subsection*{Script}
Bengali is an \textbf{abugida}: a consonant already carries a vowel, and you write a vowel mark only when you want a different one.

\begin{quote}
\bn{ন}
\end{quote}

Now the part that catches every reader who comes to Bengali from Hindi or Sanskrit. The inherent vowel here is \textbf{not} \emph{a}. It is \textbf{ɔ} — the vowel of English \emph{awe}, rounded and open. So this letter alone is \textbf{nɔ}, not \emph{na}.

Devanagari's bare consonant says \emph{ka}; Bengali's says \emph{kɔ}. Same family, same system, \textbf{different default}. It is one of the first things that makes Bengali sound like Bengali rather than like Hindi read aloud.

Look at the top of the shape: a \textbf{horizontal bar}, and in running text it joins up across a whole word so the writing hangs from a line — the same head-line Devanagari has, and the one Gujarati erases.

\subsection*{Your turn}
\begin{itemize}
\raggedright
  \item \emph{Write it:} \bn{ন}
  \item \emph{Read it:} every piece this chapter has given you so far, in order
\end{itemize}

\begin{tcolorbox}[breakable,colback=teal!4,colframe=teal!35!black,title={Before you move on}]
What vowel does a bare Bengali consonant carry? (\textbf{ɔ}, the vowel of \emph{awe} — not \emph{a}.) And what does Devanagari's carry instead? (\emph{\textbf{a}}.)

Source: \href{https://www.unicode.org/charts/PDF/U0980.pdf}{Unicode Bengali chart}.
\end{tcolorbox}
\section[-a]{\bn{া} — the vowel sign that replaces the inherent vowel}

\label{lesson:BN-W01-aa-matra}



\begin{quote}
Write the piece from the previous lesson before adding another.
\end{quote}

\subsection*{Script}
To get a different vowel you \textbf{replace} the inherent one with a mark called a \textbf{mātrā}.

\begin{quote}
\bn{া}
\end{quote}

A single upright stroke standing to the \textbf{right} of its consonant, dropping from the head-line. Attached, it swaps the inherent \textbf{ɔ} for a long open \textbf{a}.

\begin{quote}
\textbf{\bn{ন}} + \textbf{\bn{া}} $\to$ \textbf{\bn{না}}
\end{quote}

\textbf{na} — and that is the word for \textbf{no}, which you have been saying since your first lesson.

One consonant and one mark, and a word you could already say became a word you can \textbf{read}.

\subsection*{Your turn}
\begin{itemize}
\raggedright
  \item \emph{Write it:} \bn{া}
  \item \emph{Read it:} every piece this chapter has given you so far, in order
\end{itemize}

\begin{tcolorbox}[breakable,colback=teal!4,colframe=teal!35!black,title={Before you move on}]
What does a mātrā do to the inherent vowel? (\textbf{Replaces} it.) What word did those two pieces make? (\textbf{\bn{না}} — \emph{na}, \textquotedblleft{}no\textquotedblright{}.)

Source: \href{https://www.unicode.org/charts/PDF/U0980.pdf}{Unicode Bengali chart}.
\end{tcolorbox}
\section[a]{\bn{আ} — the same vowel standing alone}

\label{lesson:BN-W01-aa}



\begin{quote}
Write the piece from the previous lesson before adding another.
\end{quote}

\subsection*{Script}
The mark you just learned only works \textbf{attached to a consonant}. So what do you write when a word \textbf{begins} with that vowel?

A different character:

\begin{quote}
\bn{আ}
\end{quote}

That is the same \textbf{a} sound, but as an \textbf{independent vowel} — a full letter that stands on its own at the start of a word.

\textbf{These are two different characters, not two styles of one.} A font, a keyboard and this book all treat them as entirely separate, and confusing them is the commonest spelling error in every Indic script.

What decides which one you write is \textbf{position}, not sound:

\begin{itemize}
\raggedright
  \item \textbf{\bn{আ}} begins a word.
  \item \textbf{\bn{া}} rides on a consonant.
\end{itemize}

Two words you already say begin with this letter — the one for \emph{I see you later} and the one for \emph{all right}.

\subsection*{Your turn}
\begin{itemize}
\raggedright
  \item \emph{Write it:} \bn{আ}
  \item \emph{Read it:} every piece this chapter has given you so far, in order
\end{itemize}

\begin{tcolorbox}[breakable,colback=teal!4,colframe=teal!35!black,title={Before you move on}]
Which of the two starts a word? (\textbf{\bn{আ}}.) Which rides on a consonant? (\textbf{\bn{া}}.) Are they the same character? (\textbf{No} — two separate characters, one sound.)

Source: \href{https://www.unicode.org/charts/PDF/U0980.pdf}{Unicode Bengali chart}.
\end{tcolorbox}
\section[hɔ]{\bn{হ} — the consonant ha}

\label{lesson:BN-W01-ha}



\begin{quote}
Write the piece from the previous lesson before adding another.
\end{quote}

\subsection*{Script}
\begin{quote}
\bn{হ}
\end{quote}

\textbf{hɔ} — the \emph{h} of English \emph{hat}, plus the inherent vowel you now know is \textbf{ɔ} rather than \emph{a}.

No new machinery: every rule you hold applies unchanged. \textbf{New shapes are coming; new rules are not.}

You have been saying this letter since your first lesson — it opens the word for \emph{yes}.

\subsection*{Your turn}
\begin{itemize}
\raggedright
  \item \emph{Write it:} \bn{হ}
  \item \emph{Read it:} every piece this chapter has given you so far, in order
\end{itemize}

\begin{tcolorbox}[breakable,colback=teal!4,colframe=teal!35!black,title={Before you move on}]
What is the inherent vowel in \textbf{\bn{হ}}? (\textbf{ɔ}.) How many new rules did this lesson add? (\textbf{None}.)

Source: \href{https://www.unicode.org/charts/PDF/U0980.pdf}{Unicode Bengali chart}.
\end{tcolorbox}
\section[mɔ]{\bn{ম} — the consonant ma}

\label{lesson:BN-W01-ma}



\begin{quote}
Write the piece from the previous lesson before adding another.
\end{quote}

\subsection*{Script}
\begin{quote}
\bn{ম}
\end{quote}

\textbf{mɔ} — the \emph{m} of English \emph{met}, plus the inherent vowel.

Three consonants now. Say them in the order you met them, each with the vowel nobody wrote:

\begin{quote}
\textbf{\bn{ন}}  \textbf{\bn{হ}}  \textbf{\bn{ম}}
\end{quote}

\emph{nɔ, hɔ, mɔ.} Notice you are saying \emph{aw} three times, not \emph{ah}.

\subsection*{Your turn}
\begin{itemize}
\raggedright
  \item \emph{Write it:} \bn{ম}
  \item \emph{Read it:} every piece this chapter has given you so far, in order
\end{itemize}

\begin{tcolorbox}[breakable,colback=teal!4,colframe=teal!35!black,title={Before you move on}]
Say the three consonants you hold, each with its inherent vowel. (\emph{\textbf{nɔ, hɔ, mɔ}}.)

Source: \href{https://www.unicode.org/charts/PDF/U0980.pdf}{Unicode Bengali chart}.
\end{tcolorbox}
\section[shɔ]{\bn{স} — the consonant that is written s and said sh}

\label{lesson:BN-W01-sa}



\begin{quote}
Write the piece from the previous lesson before adding another.
\end{quote}

\subsection*{Script}
\begin{quote}
\bn{স}
\end{quote}

Here the letter and the sound part company, and it is worth knowing early.

This letter descends from the Sanskrit \textbf{s}, and every transliteration writes it \emph{s}. \textbf{In Bengali it is normally pronounced \emph{sh}} — the \emph{sh} of English \emph{shoe}.

So the word you say as \emph{nomosh-kar} is written with what looks like an \emph{s}. The spelling records the word's Sanskrit ancestry; the pronunciation records what Bengali did with it afterwards. \textbf{Both are true, and the script keeps the older one.}

English does exactly this with \emph{knight} and \emph{through}. Every literary language accumulates spellings that outlive their sounds.

\subsection*{Your turn}
\begin{itemize}
\raggedright
  \item \emph{Write it:} \bn{স}
  \item \emph{Read it:} every piece this chapter has given you so far, in order
\end{itemize}

\begin{tcolorbox}[breakable,colback=teal!4,colframe=teal!35!black,title={Before you move on}]
How is \textbf{\bn{স}} written in transliteration, and how is it said? (Written \emph{\textbf{s}}, said \emph{\textbf{sh}}.) Why the difference? (The spelling keeps the \textbf{Sanskrit ancestry}; the sound moved on.)

Source: \href{https://www.unicode.org/charts/PDF/U0980.pdf}{Unicode Bengali chart}.
\end{tcolorbox}
\section[kɔ]{\bn{ক} — the consonant ka}

\label{lesson:BN-W01-ka}



\begin{quote}
Write the piece from the previous lesson before adding another.
\end{quote}

\subsection*{Script}
\begin{quote}
\bn{ক}
\end{quote}

\textbf{kɔ} — the \emph{k} of English \emph{skip}, with \textbf{no puff of breath} after it.

Hold a palm in front of your mouth and say English \emph{kit}: you feel air. Say \emph{skip}: you do not. \textbf{This is the second one.} Bengali has a separate letter for the breathy version, and swapping them changes the word — so a difference English treats as an accident, this script writes down.

Five consonants, one mātrā, one independent vowel. One piece of machinery left.

\subsection*{Your turn}
\begin{itemize}
\raggedright
  \item \emph{Write it:} \bn{ক}
  \item \emph{Read it:} every piece this chapter has given you so far, in order
\end{itemize}

\begin{tcolorbox}[breakable,colback=teal!4,colframe=teal!35!black,title={Before you move on}]
Is \textbf{\bn{ক}} the \emph{k} of \emph{kit} or of \emph{skip}? (Of \emph{\textbf{skip}} — no puff of air.) Does the script care? (\textbf{Yes} — the breathy one is a different letter.)

Source: \href{https://www.unicode.org/charts/PDF/U0980.pdf}{Unicode Bengali chart}.
\end{tcolorbox}
\section[(hasanta)]{\bn{্} — the mark that DELETES the inherent vowel and fuses two consonants}

\label{lesson:BN-W01-hasanta}



\begin{quote}
Write the piece from the previous lesson before adding another.
\end{quote}

\subsection*{Script}
The inherent vowel is a nuisance when you do not want it. If a word presses two consonants together with no vowel between them, writing them plainly would insert an \textbf{ɔ} that is not there.

So there is a mark that removes it:

\begin{quote}
\bn{্}
\end{quote}

Bengali calls it the \textbf{hasanta} — Devanagari calls the same thing a \emph{virama}. It sits under a consonant and means: \emph{no vowel here.}

Then something happens that no European alphabet does. A consonant stripped of its vowel \textbf{fuses} with the one after it into a single joined shape, a \textbf{conjunct}:

\begin{quote}
\textbf{\bn{স}} + \textbf{\bn{্}} + \textbf{\bn{ক}} $\to$ \textbf{\bn{স্ক}}
\end{quote}

One cluster. \textbf{Conjuncts are not new characters to memorise} — they are two you already know, joined. Bengali's are more thoroughly fused than Devanagari's, so some take a moment to take apart, but the principle does not change.

\subsection*{Your turn}
\begin{itemize}
\raggedright
  \item \emph{Write it:} \bn{্}
  \item \emph{Read it:} every piece this chapter has given you so far, in order
\end{itemize}

\begin{tcolorbox}[breakable,colback=teal!4,colframe=teal!35!black,title={Before you move on}]
What does the hasanta do? (\textbf{Removes the inherent vowel}.) What is Devanagari's name for the same mark? (The \textbf{virama}.) And are conjuncts new letters? (\textbf{No} — two known letters, fused.)

Source: \href{https://www.unicode.org/charts/PDF/U0980.pdf}{Unicode Bengali chart}.
\end{tcolorbox}
\section[rɔ]{\bn{র} — the last consonant of the greeting}

\label{lesson:BN-W01-ra}



\begin{quote}
Write the piece from the previous lesson before adding another.
\end{quote}

\subsection*{Script}
\begin{quote}
\bn{র}
\end{quote}

\textbf{rɔ} — a light tap of the tongue, closer to the \emph{r} of Spanish \emph{pero} than to any English \emph{r}.

That is the last piece. Everything the greeting needs is now something you can write from memory.

\subsection*{Your turn}
\begin{itemize}
\raggedright
  \item \emph{Write it:} \bn{র}
  \item \emph{Read it:} every piece this chapter has given you so far, in order
\end{itemize}

\begin{tcolorbox}[breakable,colback=teal!4,colframe=teal!35!black,title={Before you move on}]
How is \textbf{\bn{র}} made? (A \textbf{light tap}, not an English \emph{r}.)

Source: \href{https://www.unicode.org/charts/PDF/U0980.pdf}{Unicode Bengali chart}.
\end{tcolorbox}
\section[nomoshkar]{\bn{নমস্কার} — seven pieces, all of them already yours}

\label{lesson:BN-W01-nomoshkar-read}



\begin{quote}
Write the piece from the previous lesson before adding another.
\end{quote}

\subsection*{Script}
Nothing new is introduced here. Every piece below is one you can write.

\begin{quote}
\textbf{\bn{ন}} + \textbf{\bn{ম}} + \textbf{\bn{স}} + \textbf{\bn{্}} + \textbf{\bn{ক}} + \textbf{\bn{া}} + \textbf{\bn{র}} $\to$ \textbf{\bn{নমস্কার}}
\end{quote}

Read it the way you built it:

\begin{itemize}
\raggedright
  \item \textbf{\bn{ন}} — \emph{nɔ}, inherent vowel
  \item \textbf{\bn{ম}} — \emph{mɔ}, inherent vowel
  \item \textbf{\bn{স্ক}} — the conjunct: the \emph{s}-letter with its vowel removed, fused to \emph{kɔ}
  \item \textbf{\bn{া}} — the mātrā, making that piece \emph{ka}
  \item \textbf{\bn{র}} — \emph{rɔ}
\end{itemize}

\textbf{nɔ-mosh-kar} — and notice the \emph{sh}, from a letter written \emph{s}.

Count what it took: \textbf{six consonants, one mātrā, one independent vowel and one hasanta.} With them you can read \emph{no} and the greeting, and every consonant you meet from here obeys rules you already hold.

Bengali has around forty more letters. It does not have any more \textbf{systems}.

\subsection*{Your turn}
\begin{itemize}
\raggedright
  \item \emph{Write it:} \bn{নমস্কার}
  \item \emph{Read it:} every piece this chapter has given you so far, in order
\end{itemize}

\begin{tcolorbox}[breakable,colback=teal!4,colframe=teal!35!black,title={Before you move on}]
How many new pieces were in this lesson? (\textbf{None}.) What is \textbf{\bn{স্ক}} made of? (The \textbf{s}-letter with its vowel removed, joined to \textbf{\bn{ক}}.) And what still remains to learn — rules, or shapes? (\textbf{Shapes}.)

Source: \href{https://www.unicode.org/charts/PDF/U0980.pdf}{Unicode Bengali chart}.
\end{tcolorbox}
